\documentclass{leveldoc}
\DocUses{foundations,first_bend,slingshot}
\DocFingerprint{7eb591287c464a50}

\begin{document}

\section{What the level asks}\label{sec:asks}

The game's text:
\leveldescription

Two electrons ($q = \val{q}$) start together at $(\val{xZero}, \val{yZero})$, both moving along
$+x$, at two energies: $T_1 = \val{TOne}$ ($v = \val{betaOne}\,c$) and $T_2 = \val{TTwo}$
($v = \val{betaTwo}\,c$). Each must reach its own detector on the left of the board: electron 1
the rectangle $\val{xDetLoOne} \le x \le \val{xDetHiOne}$,
$\val{yDetLoOne} \le y \le \val{yDetHiOne}$, electron 2 the rectangle
$\val{xDetLoTwo} \le x \le \val{xDetHiTwo}$, $\val{yDetLoTwo} \le y \le \val{yDetHiTwo}$ (no
acceptance windows, no kinetic windows, no gates). The player places up to three charges from
the level's magnitudes in the region $\val{regX0} \le x \le \val{regX1}$,
$\val{regY0} \le y \le \val{regY1}$. A flight ends when the electron enters its detector,
touches a charge, or leaves the world (\xsecref{foundations}{sec:board}).

The reference places one positive charge $Q = \val{Q}$ at $(\val{xC}, \val{yC})$, straight below
the launch. Its attraction bends both electrons round it: electron 1 exactly on a circle, which
meets its detector after a little more than half a turn; electron 2, too slow for a circle,
spirals in and meets its own detector on the way. The radius of a circular orbit is set by the
relativistic $pv$, not by Newton's $2T$: the level's lesson. With one charge the flights are
Level 1's closed form (\levelref{first_bend}); this document applies it, and checks it against a
numerical march.

\section{Setup}\label{sec:setup}

The level file, field by field (the appendix reproduces it):
\setuptable

In symbols: both electrons have charge $q = \val{q}$, mass $m = \val{m}$ and radius $0$; the
speed of light is $c = \val{c}$. The coupling with the reference charge is
$\kappa = qQ = \val{kappa}$: attractive. The launch is $r_0 = \val{rZero}$ above the charge,
moving perpendicular to the radius, so it is an apsis of both orbits. By
\xeqref{foundations}{eq:launch}, electron 1 has $\gamma = \val{gammaOne}$ and
$p = \val{pOne}$, electron 2 $\gamma = \val{gammaTwo}$ and $p = \val{pTwo}$. The world is
$-1 \le x \le \val{xWorld}$, $-1 \le y \le \val{yWorld}$.

\begin{figure}[htbp]
\centering
\begin{tikzpicture}
\begin{axis}[axis equal image, width=0.92\linewidth, xmin=-1.6, xmax=31.6, ymin=-1.6, ymax=21.6,
  xtick={0,5,...,30}, ytick={0,5,...,20}, tick label style={font=\small},
  xlabel={$x$}, ylabel={$y$}, label style={font=\small}]
  \draw[black!8] (0,0) grid (30,20);
  \draw[black!45, dashed] (-1,-1) rectangle (\rawval{xWorld},\rawval{yWorld});
  \draw[black!30, dotted] (\rawval{regX0},\rawval{regY0}) rectangle (\rawval{regX1},\rawval{regY1});
  \fill[green!50!black, fill opacity=0.18] (\rawval{xDetLoOne},\rawval{yDetLoOne}) rectangle (\rawval{xDetHiOne},\rawval{yDetHiOne});
  \draw[green!40!black] (\rawval{xDetLoOne},\rawval{yDetLoOne}) rectangle (\rawval{xDetHiOne},\rawval{yDetHiOne});
  \fill[green!50!black, fill opacity=0.18] (\rawval{xDetLoTwo},\rawval{yDetLoTwo}) rectangle (\rawval{xDetHiTwo},\rawval{yDetHiTwo});
  \draw[green!40!black] (\rawval{xDetLoTwo},\rawval{yDetLoTwo}) rectangle (\rawval{xDetHiTwo},\rawval{yDetHiTwo});
  \addplot[very thick, blue!70!black] table {pathone.dat};
  \addplot[very thick, red!70!black] table {pathtwo.dat};
  \fill[red!75!black] (\rawval{xC},\rawval{yC}) circle[radius=\rawval{R}];
  \fill (\rawval{xZero},\rawval{yZero}) circle[radius=0.15];
\end{axis}
\end{tikzpicture}
\caption{The reference charge (red disc) and the two electrons' flights: electron 1 (blue) on a
circle of radius $r_0$, entering its detector (lower green box) through its right face; electron
2 (red) spiralling in and entering its detector (upper green box) through its bottom face.}
\label{fig:board}
\end{figure}

\section{Principles and model}\label{sec:model}

\begin{principles}
  \item The field of a point charge and the force on a non-polarizable particle,
    \xeqref{foundations}{eq:coulomb} and \xeqref{foundations}{eq:sphereforce}; the relativistic
    equation of motion \xeqref{foundations}{eq:eom}, the launch \xeqref{foundations}{eq:launch},
    the energy \xeqref{foundations}{eq:energy} and the angular momentum
    \xeqref{foundations}{eq:angmom}.
  \item The relativistic orbit about one charge, \xeqref{first_bend}{eq:orbit} and its solution
    \xeqref{first_bend}{eq:general}, with the time \xeqref{first_bend}{eq:dtdth} and the momentum
    \xeqref{first_bend}{eq:momentum}.
  \item The board's rules; the classical Runge--Kutta method, cubic Hermite interpolation, the
    error estimate by halving the step, Simpson's rule (\xsecref{foundations}{sec:board},
    \xsecref{foundations}{sec:numerics}), as in \xsecref{slingshot}{sec:method}.
\end{principles}

The assumptions are Level 1's: the charge is held in place, the electron is a point, and its
radiation is not taken from it (\xsecref{first_bend}{sec:model}; the game computes the
radiated energy and checks it below $10^{-10}$ of the launch energy).

\section{Derivation}\label{sec:derivation}

\subsection{The circular orbit}\label{sec:circle}

A circle of radius $r$ about the charge needs the attraction to supply the centripetal force:
$\abs{\kappa}/r^2 = \abs{\dot{\vect{p}}} = p\,\dot\varphi = pv/r$, because on a circle
$\abs{\vect{p}}$ is constant and turns at the rate $v/r$. So
\begin{equation}\label{eq:circle}
  \abs{\kappa} = p\,v\,r = \gamma m v^2 r .
\end{equation}
Newton's mechanics would ask $\abs{\kappa} = mv^2r = 2Tr$. For electron 1, $pv\,r_0 = \val{pvrOne}$
equals $\abs{\kappa}$: its launch is exactly circular, while Newton's condition would need
$2T_1r_0 = \val{newtonOne}$. For electron 2, $pv\,r_0 = \val{pvrTwo}$, less than $\abs{\kappa}$:
the attraction wins, and it falls inwards.

\subsection{The orbits from Level 1's closed form}\label{sec:orbits}

With $\theta$ the angle swept since the launch and $u = 1/r$, the orbit obeys
\xeqref{first_bend}{eq:orbit}, $u'' + \Gamma^2u = -\kappa W/(c^2L^2)$ with
$\Gamma^2 = 1 - \kappa^2/(c^2L^2)$, $W$ the energy and $L = r_0p$ the angular momentum. The
launch is an apsis, $u'_0 = 0$, so \xeqref{first_bend}{eq:general} gives
\begin{equation}\label{eq:spiral}
  u(\theta) = B + (u_0 - B)\,C(\theta), \qquad B = -\frac{\kappa W}{c^2L^2\Gamma^2},
  \qquad C(\theta) = \begin{cases} \cos\Gamma\theta, & \Gamma^2 > 0,\\
    \cosh\abs{\Gamma}\theta, & \Gamma^2 < 0. \end{cases}
\end{equation}
Electron 1: $W = \val{WOne}$, $L = \val{LOne}$, $\Gamma^2 = \val{GammaSqOne}$, so
$\Gamma = \val{GammaOne}$, and $B = \val{BOne}$, equal to $u_0 = 1/r_0$: the amplitude
$u_0 - B = \val{AOne}$ vanishes to rounding and the orbit is the circle $r = r_0$. On a circle
\eqref{eq:circle} gives $\kappa^2 = (pvr)^2 = c^2L^2v^2/c^2$, so $\Gamma^2 = 1 - v^2/c^2$ and
$\Gamma = 1/\gamma$, as found.

Electron 2: $W = \gamma_0mc^2 + \kappa/r_0 = \val{WTwo}$ exactly ($\gamma_0 mc^2 = \abs{\kappa}/r_0$
here), so $B = 0$; $L = \val{LTwo}$ and $\Gamma^2 = \val{GammaSqTwo}$, negative: the
hyperbolic case,
\[
  r(\theta) = \frac{r_0}{\cosh\abs{\Gamma}\theta}, \qquad \abs{\Gamma} = \val{GammaTwo},
\]
a spiral that reaches the charge after infinitely many turns in $\theta$ (in a finite time).
It passes the detector long before.

\subsection{The entry, the time and the momentum}\label{sec:entry}

The position is $\vect{r}_\mathrm{c} + r(\theta)(\cos\varphi, \sin\varphi)$ with
$\varphi = \varphi_0 - \theta$ (both turn clockwise). The entry is the first $\theta$ at which
it lies in the detector, found by steps of $10^{-3}$ in $\theta$ and bisection; the charge's
disc and the world's edge are checked on the same steps and are not reached first. The time is
\xeqref{first_bend}{eq:dtdth} integrated by Simpson's rule with $\val{nSimpson}$ panels (its
error estimated by halving them); on the circle the integrand is constant and Simpson's rule
exact. The momentum is \xeqref{first_bend}{eq:momentum}.

\subsection{An independent check}\label{sec:march}

The equations of motion are also marched by the classical Runge--Kutta method
\xeqref{foundations}{eq:rk4} at $h = \val{H}$, the entry located on the cubic Hermite
interpolant (\xsecref{slingshot}{sec:method}). Its end states agree with the closed form to
$\val{marchOne}$ (electron 1) and $\val{marchTwo}$ (electron 2); its own step-halving estimates
are $\val{mErrOne}$ and $\val{mErrTwo}$, and the energy drifts by at most $\val{wDevTwo}$.

\section{The reference solution by hand}\label{sec:hand}

\begin{procedure}
  \item\label{st:launch} \emph{The launches.} Electron 1: $\gamma = \val{gammaOne}$,
    $p = \val{pOne}$; electron 2: $\gamma = \val{gammaTwo}$, $p = \val{pTwo}$; both at
    $r_0 = \val{rZero}$, an apsis.
  \item\label{st:circle} \emph{Electron 1, the circle.} $pv\,r_0 = \val{pvrOne} = \abs{\kappa}$:
    $r = r_0$ throughout. It reaches its detector's right face, $x = \val{xDetHiOne}$, after
    $\theta = \val{thEdegOne}$ degrees, at $(\val{xEOne}, \val{yEOne})$, at
    $t = \val{tEOne}$ (Simpson's error estimate $\val{tErrOne}$), with momentum
    $(\val{pxEOne}, \val{pyEOne})$.
  \item\label{st:spiral} \emph{Electron 2, the spiral.} $r = r_0/\cosh(\val{GammaTwo}\,\theta)$.
    It reaches its detector's \val{faceTwo} face, $y = \val{yDetLoTwo}$, after
    $\theta = \val{thEdegTwo}$ degrees, at $(\val{xETwo}, \val{yETwo})$, at $t = \val{tETwo}$
    (Simpson's error estimate $\val{tErrTwo}$), with momentum $(\val{pxETwo}, \val{pyETwo})$:
    much faster than at the launch, having fallen towards the charge.
\end{procedure}

\begin{result}[The reference flights]
Electron 1 enters its detector at $(\val{xEOne}, \val{yEOne})$ at $t = \val{tEOne}$; electron 2
enters its detector at $(\val{xETwo}, \val{yETwo})$ at $t = \val{tETwo}$. The level is solved.
\end{result}

\section{Comparison with the game}\label{sec:comparison}

\comparisontable

All differences are a few $10^{-11}$ or less: the game's verified accuracy (about $10^{-10}$,
Foundations §6), against our estimates of a few $10^{-13}$. The tolerance is $10^{-9}$, ten
times the larger of the two computations' estimates, or the game's own gap between its preview
and its verification where that is larger (electron 2, which falls into a strong field:
$\val{gapTwo}$), fixed before comparing.

\section{Other solutions}\label{sec:others}

On the empty board, the second placement in \code{placements.json}, both electrons fly straight
through the \val{bareFace} edge at $(\val{bareX}, \val{bareY})$. One charge straight below the
launch, at the distance that makes electron 1's orbit circular by \eqref{eq:circle}, is the
reference's idea; the detectors are placed round its flights. Other placements bend the two
electrons differently; whether a single charge elsewhere serves both is a question for the
game's verdicts, which this document does not map.

\section{What the game computes}\label{sec:engine}

The game integrates the same equations of motion with its adaptive eighth-order Runge--Kutta
method (\physref{5}), locating the events by its dense output (\physref{6}) and verifying the
outcome at a tighter tolerance (\physref{7}). It uses no closed form; the agreement in the
comparison is that of the integration with Level 1's exact orbit.

\section{Further reading}\label{sec:reading}

The relativistic Kepler problem and its precession (the factor $\Gamma$), and the capture of a
particle whose angular momentum falls below $\abs{\kappa}/c$: \LandauII{§39}; relativistic
mechanics: \Goldstein{Ch.~7}. The closed form used here: \levelref{first_bend}.

\levelfileappendix

\end{document}
